\scanpage{537}%
\mspage{448}%
% archive sheet number: 536
la forme initiale $2^\circ)$ de la condition 2bis que l'on a $\tau_{j''j}=p$ pour tt $j''\geq j'$, donc si $j$ n'est pas objet final de $J$, OPS $j'>j$. Mais $J$ n'a pas d'objet final, car sinon $I$ en aurait un, cas que l'on pouvait éliminer.) Donc il existe $j'$ tel que $j'\overset K>j$. Or on a ceci (qu'on appliquera au cas $i\in K'$)

\[i\leq j\quad\text{et}\quad j\overset K<j'\ \Longrightarrow\ i\overset K<j'.\]

\diagram{\begin{tikzcd}[ampersand replacement=\&] x_i \arrow[d,equal]\arrow[r,shift left,"{\tau_{j,i}\deleted{\unsure{*}}}"] \& x_j\arrow[l,shift left,"{\unsure{\leftarrow|\deleted{\leftarrow}}}" description,draw=none]\arrow[r,"{\deleted{\unsure{\tau_{j'j}}}\inserted{\tau_{j',j}=p}}"] \& x_{j'} \\ x_0 \& \& \end{tikzcd}}

\[\begin{aligned}\tau_{j'i}&=\tau_{j'j}\tau_{ji}={}&p\tau_{ji}&\overset{?}{=}p\\ &&\underset{\tau_{ji}p}{\Vert}\ \text{par }1^\circ\end{aligned}\]

\[\text{on a en effet}\quad\left.\begin{aligned}\tau_{j',i}&=p\,\tau_{j'i}\\ \tau^{\unsure{\prime|{}}}_{j',i}&=\tau_{j'i}p\end{aligned}\right|\Longrightarrow\underset{\text{par }4.15.7.4}{\tau_{j'i}=p}.\]

Cela achève la démonstration du corollaire 4.15.7.3

\inserted{\uline{Cor.} 4.15.7.5 \inserted{(de 4.15.7.3)} Soit $(x_i)_{i\in I}$ un syst. inductif \inserted{ess. constant,} $I$ grd devant $\pi_0\geq\operatorname{card} C$, dans $C$ \uline{sans projecteurs}. Alors il existe $J\to I$ cofinal, avec $J$ ordonné grd devant $\pi_0$, et tel que le syst. inductif induit sur $J$ soit \uline{constant}.}

\uline{Démontrons le lemme} 4.15.7.1. Il faut expliciter l'existence de deux flèches dans $\mathrm{Ind}(C)$

\diagram{\begin{tikzcd}[ampersand replacement=\&] z \arrow[r,shift left,"\alpha","{\unsure{\deleted{\leftarrow}|{}}}"{description,pos=0.15}] \& \mathop{\text{\textquotedblleft}\lim\text{\textquotedblright}}\limits_{i\in I}x_i=\underline{x}\arrow[l,shift left,"\beta"] \end{tikzcd}}

avec $\beta\alpha=\mathrm{id}_z$, $\alpha\beta=\mathrm{id}_{\underline{x}}$. \deleted{\unsure{La}} La donnée
